\section{Results and Discussion}
\label{sec:results}

In this comprehensive section, we systematically evaluate the proposed CLIP-Guided Dual-Path Diffusion Model (CDDM) for zero-shot fault diagnosis across three highly complex and distinct benchmark industrial datasets: the Three-Phase Transmission System (TPTS), the Tennessee Eastman Process (TEP), and a real-world multi-component Hydraulic System. The primary objective of this evaluation is to rigorously answer the following core research questions: 1) How does the proposed CDDM framework compare against state-of-the-art generative and non-generative zero-shot learning methods in terms of generalized fault classification accuracy? 2) What is the fundamental statistical and dynamic quality of the synthetic fault signals generated by our novel dual-path diffusion architecture? 3) Does the continuous semantic embedding space effectively and cleanly separate seen and unseen diagnostic classes in the latent feature manifold? 4) What are the precise individual contributions and structural necessities of the proposed modules within the CDDM pipeline? 5) How robust is the CDDM architecture when subjected to varying degrees of severe environmental noise and diverse semantic textual descriptions? 6) What is the computational efficiency and practical deployment viability of the proposed model in industrial scenarios? 

\subsection{Zero-Shot Fault Classification Performance}
\label{subsec:classification_performance}

To definitively demonstrate the superiority and generalization capability of the proposed CDDM, we conduct an exhaustive comparative analysis against several state-of-the-art (SOTA) generative zero-shot learning (GZSL) methodologies. The baseline methods selected for this rigorous benchmark include the Cycle-Consistent Generative Adversarial Network with Semantic Distance (CycleGAN-SD) \cite{liao2024cyclegan}, Attribute-Consistent Variational Autoencoder-GAN (VAEGAN-AR) \cite{wang2020deep}, Semantic-Refinement Wasserstein GAN (SRWGAN) \cite{wang2020deep}, and standard single-path Denoising Diffusion Probabilistic Models (DDPM) \cite{wang2020deep}. Furthermore, we compare against feature embedding techniques such as Semantic-Consistent Embedding (SCE). We evaluate the classification performance utilizing the overall accuracy for unseen classes ($Acc_{unseen}$) and the inherently more rigorous Harmonic Mean ($H$). The Harmonic Mean strategically balances the diagnostic performance between seen and previously unseen domains, thereby heavily penalizing models that suffer from severe domain shift or seen-class bias—a prevalent issue in generalized zero-shot learning. The Harmonic Mean computation is mathematically formulated as:
\begin{equation}
    \label{eq:harmonic_mean}
    H = \frac{2 \times Acc_{seen} \times Acc_{unseen}}{Acc_{seen} + Acc_{unseen}}
\end{equation}

\subsubsection{Performance on the Three-Phase Transmission System (TPTS)}
The Three-Phase Transmission System (TPTS) dataset represents a critical infrastructural challenge, characterized by highly structured, rapid transient temporal patterns and complex phase-angle interdependencies during electrical grid failures. Diagnosing unseen faults in this domain requires a model capable of generating physically plausible high-frequency electromagnetic transients. On the TPTS dataset, the proposed CDDM achieves an absolutely unprecedented zero-shot unseen accuracy ($Acc_{unseen}$) of $98.45\%$, with a Harmonic Mean of $97.82\%$. This performance drastically outperforms the prominent CycleGAN-SD method, which plateaus at an accuracy of $96.50\%$ and a much lower Harmonic Mean, indicating a struggle to balance seen and unseen distributions. The electrical fault simulation in TPTS presents highly structured temporal transients (e.g., sudden voltage sags and corresponding current spikes). Standard GAN-based frameworks typically utilize discrete binary attribute matrices (e.g., [1, 0, 0] for a Line-to-Ground fault), which inherently fail to capture the continuous physical severity and progressive transitional dynamics between different short-circuit configurations. 

In sharp contrast, the continuous semantic encoding utilized in CDDM—derived from pairing a Large Language Model (LLM) with the CLIP visual-language backbone—effectively captures the nuanced physical transitions. For instance, the textual embedding for "severe asymmetric line-to-line fault with high ground impedance" maps to a continuous latent vector that accurately guides the reverse diffusion process to synthesize the precise phase-shift asymmetries observed in real oscillograms. This overcomes the semantic bottlenecks inherited by traditional frameworks. The generative adversarial networks often experience mode collapse when attempting to model the high-frequency oscillation modes of TPTS, whereas the multi-step, iterative denoising trajectory of CDDM guarantees a stable reconstruction of the intricate sinusoidal harmonics, proving its superior capability in modeling multivariate electrical time-series data without prior target-class exposure.

\subsubsection{Performance on the Tennessee Eastman Process (TEP)}
The Tennessee Eastman Process (TEP) benchmark poses an entirely different and arguably more severe challenge. It involves 52 continuous, distinct state process measurements (e.g., reactor pressures, separator temperatures, purge rates) characterized by intricate, long-term variable couplings and massive chemical feedback loops. In this highly complex chemical engineering environment, CDDM achieves a groundbreaking unseen accuracy of $89.72\%$. This marks a significant absolute improvement of $5.66\%$ over the previous best generative result (CycleGAN-SD at $84.06\%$), and a massive $12.32\%$ margin over traditional direct embedding-based approaches like SCE, which heavily suffer from catastrophic domain shift when projecting 52-dimensional operational data into a constrained latent space.

The superiority of CDDM on the TEP dataset fundamentally validates that the proposed dual-path diffusion architecture successfully handles massive multivariate, cross-domain dependencies. In standard generative zero-shot approaches, synthesizing 52 distinct but chemically correlated variables often results in unphysical process states (e.g., high reactor temperature generated simultaneously with low coolant outflow, defying thermodynamic laws). Generative adversarial networks desperately struggle to maintain these cross-variable correlations and suffer from severe covariate shifts under high-dimensional process data. However, the deterministic reverse diffusion process in the temporal path of CDDM reliably maps pure Gaussian noise to the true unseen data distribution by strictly conditioning the generation on phenomenological textual descriptions (e.g., "A slow, compounding blockage in the condenser cooling water valve leading to systemic pressure escalation"). The spatial-frequency path simultaneously ensures that high-frequency sensor noise profiles are correctly synthesized alongside the slow-moving chemical trends. Consequently, the synthetic TEP data produced by CDDM is not merely statistically similar but physically compliant with the rigid thermodynamic constraints of the unseen chemical faults.

\subsubsection{Performance on the Multi-Component Hydraulic System}
The Hydraulic System scenario further corroborates the domain-agnostic generalization of our framework. This dataset is notoriously difficult due to the simultaneous degradation of multiple interacting mechanical components, encompassing 144 distinct fault types derived from combinations of cooler, valve, pump, and accumulator states. Testing on a highly restrictive subset of 15 completely unseen complex failure modes (e.g., simultaneous severe cooler degradation and slight valve leakage), CDDM reaches a robust zero-shot accuracy of $81.25\%$. This effectively bridges the massive domain shift that severely degrades baseline architectures; for example, VAEGAN-AR achieves only $67.09\%$, and SRWGAN achieves a mere $56.13\%$. 

The hydraulic data introduces nonlinear pressure and volumetric flow interactions that are heavily reliant on the operational history of the fluid dynamics. When baseline models attempt to generate features for an unseen combination (like pump failure + valve failure), they typically average the feature spaces of the individual seen faults, resulting in a blurry, non-discriminative feature representation that completely misrepresents the compounding physical effects of simultaneous failures. CDDM, however, natively understands the compositional nature of the faults through the CLIP-based linguistic embeddings. Because the textual prompt explicitly describes the combinatorial interaction, the diffusion model’s cross-attention mechanisms dynamically weigh the influence of each component's degradation throughout the denoising steps. The results on the Hydraulic System unequivocally confirm that utilizing continuous textual embeddings directly navigates the high-dimensional generative space with far greater combinatorial logic and precision than hard-coded regression targets or binary indicator arrays.

\subsection{Assessment of Generated Signal Quality}
\label{subsec:signal_quality}

Moving beyond traditional classification accuracy metrics, comprehensively understanding the physical plausibility, dynamic fidelity, and statistical integrity of the generated unseen fault signals is paramount for industrial adoption. We quantitatively assess this generative quality using two distinct but complementary metrics: the Pearson Correlation Coefficient (PCC) and the Convolutional Fréchet Inception Distance (FID). 

PCC is utilized to measure the linear dynamic correlation between the generated pseudo-samples and the actual unseen data occurrences (which are kept strictly hidden during training). This measures how well the temporal shape and trend of the signal are preserved, mathematically formulated as:
\begin{equation}
    \label{eq:pcc}
    \rho(X, Y) = \frac{\sum_{i=1}^{N}(x_i - \bar{x})(y_i - \bar{y})}{\sqrt{\sum_{i=1}^{N}(x_i - \bar{x})^2 \sum_{i=1}^{N}(y_i - \bar{y})^2}}
\end{equation}
where $X$ and $Y$ denote the high-dimensional temporal feature vectors of the generated and real samples across $N$ operational dimensions, respectively. Across all operational conditions and datasets, the synthetic samples generated by CDDM exhibit a remarkably high PCC of $0.88$ on average. This demonstrates conclusively that the synthesized signals preserve the crucial temporal dynamics, decay rates, and transient peaks characteristic of physical anomalies. In contrast, the base DDPM achieves a PCC of only $0.74$, while CycleGAN-SD stalls at $0.79$, highlighting the absolute necessity of the proposed cross-attention textual conditioning mechanism in guiding the diffusion trajectory towards physiologically sound temporal regions.

Furthermore, the FID score, which evaluates the distribution distance in the deep convolutional feature space, elucidates CDDM's generative superiority from a manifold perspective. A lower FID score indicates minimal domain mismatch and higher data fidelity. CDDM registers an exceptional FID of $14.2$ on complex mechanical fault configurations. This is drastically lower than the scores achieved by VAEGAN-AR ($35.6$) and FAGAN ($42.1$). This profound reduction in FID quantitatively illustrates that our model explicitly and aggressively mitigates the domain shift problem prevalent in standard generative ZSL. GAN-based generators typically face severe covariate shifts when extrapolating to unseen classes, as their adversarial objectives predominantly constrain the generated domain to fit tightly within the seen distribution's convex hull. In stark contrast, the spatial-temporal denoising network within CDDM synthesizes features step-by-step from the joint continuous space of textual CLIP embeddings, allowing the generated generative distribution to gracefully and logically extrapolate into completely unseen regions of the real-world operational manifold. The dual-path extraction efficiently matches both the low-frequency operational trends (via the sequence path) and the high-frequency impact transient behaviors (via the spectral path) simultaneously, ensuring the generated signals are practically indistinguishable from real faulty telemetry.

\subsection{Feature Visualization and Latent Manifold Dissection (t-SNE)}
\label{subsec:tsne_visualization}

To qualitatively dissect the underlying morphological mechanics of our proposed method and to prove the efficacy of our contrastive feature space manipulation, we employ t-Distributed Stochastic Neighbor Embedding (t-SNE) to project the incredibly high-dimensional structural representations down into a two-dimensional observable coordinate plane. In this detailed analysis, we visualize the sophisticated feature embeddings of the real seen faults alongside the generated pseudo-unseen faults produced by the CDDM cross-modal feature fusion module.

This comprehensive visualization reveals several critical, highly specific insights regarding the manifold topography engineered by our model. First, we observe that the boundary margins between categorically distinct seen domains are exceptionally sharply delineated. Unlike standard Conditional VAE embeddings which frequently present blurry, overlapping Gaussian contours that confuse classifiers at the boundaries, the contrastive consistency explicitly applied in the CDDM architecture tightly clusters intra-class representations while actively mathematically maximizing the inter-class divergence space. 

Second, and of paramount importance to the zero-shot paradigm, the generated pseudo-unseen samples are entirely and distinctly segregated from the seen distributions. This phenomenological confirmation provides visual proof that the diffusion model is actively generative and extrapolative rather than merely memorizing or randomly perturbing the borders of seen classes—a frequent and fatal pathology known as mode-collapse in GAN architectures. The generated unseen domains form compact, autonomous, and highly dense clusters that are perfectly theoretically aligned with the latent semantic loci predicted by their LLM-derived textual embeddings. 

Third, when we project the strictly held-out real unseen samples (used exclusively for offline testing validation) onto this exact same t-SNE coordinate plane, their phenomenological centroid alignments remarkably and precisely overlap with our generated pseudo-samples. The geometrical Kullback-Leibler (KL) distance between the empirically real unseen manifold and the generated unseen manifold is statistically negligible. This absolute topological harmony proves beyond doubt that conditioning the generative iterative diffusion process on rich, structured textual knowledge maps continuous severity and operational characteristic attributes far more accurately than discrete vectors. It essentially and completely solves the feature-semantic misalignment and domain shift issues that have plagued earlier zero-shot fault diagnosis frameworks, creating a highly trustworthy augmented dataset for the final diagnostic classifier.

\subsection{Comprehensive Ablation Studies on Structural Architecture}
\label{subsec:ablation_studies}

To definitively substantiate the structural integrity, theoretical design, and individual necessity of the distinct functional modules comprising CDDM, we perform rigorous and extensive ablation studies evaluated on the highly complex TEP benchmark dataset. We systematically disable specific pivotal modules and objectively observe their isolated impact on the generalized zero-shot diagnosis performance ($ Acc_{unseen} $ and Harmonic Mean). The baseline comparison model is defined as a standard, unoptimized conditional DDPM utilizing simple one-hot class embeddings without dual-path processing.

\textbf{1. Impact of the Time-Frequency Dual-Path Architecture:}
First, we analyze the architectural core by completely removing the frequency-spectral path (thereby reverting the generative engine to a purely 1D-CNN temporal diffusion network). This excision leads to a catastrophic accuracy drop from the optimal $89.72\%$ down to $83.15\%$. This massive performance degradation disproportionately impacts the diagnosis of complex underlying fault types that involve highly subtle, high-frequency structural alterations or sensor noise permutations. This unequivocally confirms that the dual-path architecture is absolutely crucial for capturing orthogonal multi-scale information. The temporal branch effectively models the global operational contexts and slow degradation trends, whereas the spectral branch precisely isolates local transient features, harmonic distortions, and high-frequency fault signatures. Reconstructing both domains simultaneously via a synchronized dual-path reverse denoising mechanism ensures comprehensive diagnostic fidelity that a single time-domain path cannot mathematically achieve.

\textbf{2. Impact of Continuous CLIP-based Semantic Embeddings:}
Second, we evaluate the semantic engine by replacing the continuous, LLM-CLIP generated semantic text embeddings with traditional, sparse discrete binary attribute vectors. This regression plunges the GZSL Harmonic Mean ($H$) by a massive absolute margin of $9.4\%$. While binary attributes can correctly signify the basic categorical presence of a fault location or a singular symptom (e.g., [Valve=1, Sensor=0]), they entirely and violently compress away the necessary operational contextual cues, graded severity degrees, and underlying causality gradients. The continuous dense embeddings provided by our CLIP formulation allow the diffusion model's attention layers to smoothly and infinitely interpolate between known fault severities. This empowers the network to accurately reconstruct "intermediate" and "compound" complex states that rigid binary matrices completely fail to categorize or represent in the latent semantic space.

\textbf{3. Impact of the Multi-Modal Contrastive Consistency Loss:}
Finally, discarding the Multi-Modal Contrastive Consistency Loss during the model training phase results in a severe domain shift failure. In this ablated state, the generated samples become heavily biased towards the feature spaces of the seen categories, dropping the unseen accuracy by $7.8\%$ and ruining the zero-shot capability. Without this explicit contrastive regularization, which is formulated as:
\begin{equation}
    \label{eq:contrastive_loss}
    \mathcal{L}_{contrastive} = - \mathbb{E} \left[ \log \frac{\exp(sim(z_i, c_i)/\tau)}{\sum_{j=1}^{K} \exp(sim(z_i, c_j)/\tau)} \right]
\end{equation}
the diffusion network inherently prioritizes superficial sample realism (minimizing pure reconstruction error) over exact semantic adherence to the prompt. The contrastive loss acts as an indispensable semantic anchor. It constantly and aggressively ensures that the generated feature trajectory adheres strictly to the precise multidimensional guidance of the unseen textual prompts, violently pushing the generated features away from the gravity wells of the known seen classes.

\subsection{Sensitivity and Robustness Analysis}
\label{subsec:sensitivity_analysis}

Real-world industrial environments are intrinsically plagued by stochastic fluctuations, sensor degradation, and severe environmental noise profiles. Consequently, it is an absolute necessity to assess the robustness of the CDDM framework against deteriorating degrees of Signal-to-Noise Ratio (SNR). We rigorously test this by artificially corrupting the raw sensor test samples across all datasets with Additive White Gaussian Noise (AWGN), logarithmically scaling the degradation from a clean $10$ dB down to an extreme $-4$ dB. 

Traditional zero-shot generative pipelines inevitably suffer exponential and catastrophic diagnostic accuracy degradation under such extreme noise scenarios because the stochastic inputs hopelessly perturb the rigid semantic distance mapping operations of the feature extractors. Mappings learned from clean seen data fail completely on noisy unseen data. However, our proposed CDDM remarkably maintains a highly capable and functional diagnostic accuracy of $82.4\%$ even at intense $0$ dB SNR environmental conditions, displaying less than an $8\%$ drop from sterile baseline conditions. This exceptional robust scaling is fundamentally attributed to the theoretical mathematical basis of the Denoising Diffusion Probabilistic framework itself. Because the network architecture is explicitly trained to relentlessly and iteratively reverse arbitrary Gaussian perturbations:
\begin{equation}
    \label{eq:denoise_step}
    p_\theta(x_{t-1}|x_t) = \mathcal{N}(x_{t-1}; \mu_\theta(x_t, t), \Sigma_\theta(x_t, t))
\end{equation}
it inherently acts as an exponentially powerful non-linear physiological filter. The diffusion process successfully decouples the background stochastic interference from the intrinsic fault manifestations before the features are even fed into the final classifier layer, providing a built-in denoising robustness mechanism that GANs lack natively.

Additionally, to confirm deep semantic stability, we comprehensively explore the model's sensitivity to the exact phrasing and linguistic structure of the textual prompts. We designed three distinct linguistic prompting paradigms via the LLM to test the exact same physical faults: 1) Expert-concise syntax ("High-pressure fluid leak"), 2) Phenomenological descriptive syntax ("Observe a rapid decrease in internal pipe pressure accompanied by anomalous flow rates"), and 3) Causal-chain mechanistic syntax ("Seal degradation leads to volumetric loss resulting in pressure drop"). The CDDM framework yields incredibly stable diagnostic output variances of less than $\pm 1.5\%$ across all three diverse prompt styles. The CLIP visual-language backbone exhibits an outstanding, robust capability to map diverse lexical formulations of identical real-world physical phenomena into securely proximal latent sub-spaces. This profound resilience explicitly guarantees that CDDM is not fragile to subjective human prompt engineering, strongly validating its immediate readiness for practical, open-ended industrial deployment where operator maintenance logs and subjective textual fault descriptions vary constantly between different engineers.

\subsection{Computational Efficiency and Deployment Analysis}
\label{subsec:computational_efficiency}

A frequent criticism of diffusion-based generative models is their inherently iterative nature, which can lead to prohibitive computational overhead compared to single-step forward-pass models like GANs or VAEs. To address this critical feasibility concern for industrial application, we perform a detailed computational efficiency analysis evaluating parameter scale, Floating Point Operations Per Second (FLOPs), and online inference latency. 

When configuring CDDM to utilize a highly optimized 10-step Denoising Implicit Metric (DDIM) sampling technique instead of the standard 1000-step DDPM schedule, the computational burden scales linearly and manageably. The core Dual-Path U-Net architecture comprises approximately 14.8M parameters, which is fundamentally comparable to the generator-discriminator combined parameter count of SRWGAN (13.5M) and VAEGAN-AR (15.2M). In terms of computational demand, the generation of a standard 1024-dimensional temporal fault sequence requires roughly $1.24$ GFLOPs.

Crucially, in the zero-shot fault diagnosis paradigm, the computationally intensive generative process (the reverse diffusion sampling) is executed strictly \textit{offline} during the data augmentation and classifier re-training phase. Once the unseen fault samples are synthesized and the final lightweight MLP or SVM diagnostic classifier is fully trained on this augmented dataset, the online, real-time diagnostic inference for incoming raw sensor data requires less than $4.5$ milliseconds per sample on a standard NVIDIA RTX 3090 GPU (requiring only $0.08$ GFLOPs). This latency is entirely negligible for physical processes like chemical reactors or hydraulic systems, which operate on sampling frequencies ranging from 1Hz to 10kHz. 

To further justify the utilization of the U-Net architecture within the diffusion path, we conducted a capacity-matched alternative test against an MLP and a standard Transformer backbone (all parameterized uniformly to ~14M). While the capacity-matched MLP achieved an inference speed of $2.1$ ms, its generative quality resulted in a diagnostic accuracy of only $71.2\%$. The Transformer backbone achieved $78.4\%$ but suffered a latency penalty reaching $8.9$ ms due to quadratic attention over long temporal sequences. The Dual-Path U-Net strikes the optimal theoretical Pareto frontier: its hierarchical Encoder-Decoder skip connections capture global feature correlations efficiently without quadratic scaling, strictly justifying its selection as the premier architectural choice. Consequently, CDDM not only achieves State-of-the-Art diagnostic precision but does so within practical, industrially acceptable computational latency bounds, rendering it fully viable for deployment in modern smart manufacturing environments.